% Options for packages loaded elsewhere
\PassOptionsToPackage{unicode}{hyperref}
\PassOptionsToPackage{hyphens}{url}
\documentclass[
  11pt,
]{article}
\usepackage{xcolor}
\usepackage[margin=1in]{geometry}
\usepackage{amsmath,amssymb}
\setcounter{secnumdepth}{-\maxdimen} % remove section numbering
\usepackage{iftex}
\ifPDFTeX
  \usepackage[T1]{fontenc}
  \usepackage[utf8]{inputenc}
  \usepackage{textcomp} % provide euro and other symbols
\else % if luatex or xetex
  \usepackage{unicode-math} % this also loads fontspec
  \defaultfontfeatures{Scale=MatchLowercase}
  \defaultfontfeatures[\rmfamily]{Ligatures=TeX,Scale=1}
\fi
\usepackage{lmodern}
\ifPDFTeX\else
  % xetex/luatex font selection
\fi
% Use upquote if available, for straight quotes in verbatim environments
\IfFileExists{upquote.sty}{\usepackage{upquote}}{}
\IfFileExists{microtype.sty}{% use microtype if available
  \usepackage[]{microtype}
  \UseMicrotypeSet[protrusion]{basicmath} % disable protrusion for tt fonts
}{}
\makeatletter
\@ifundefined{KOMAClassName}{% if non-KOMA class
  \IfFileExists{parskip.sty}{%
    \usepackage{parskip}
  }{% else
    \setlength{\parindent}{0pt}
    \setlength{\parskip}{6pt plus 2pt minus 1pt}}
}{% if KOMA class
  \KOMAoptions{parskip=half}}
\makeatother
\setlength{\emergencystretch}{3em} % prevent overfull lines
\providecommand{\tightlist}{%
  \setlength{\itemsep}{0pt}\setlength{\parskip}{0pt}}
\usepackage{bookmark}
\IfFileExists{xurl.sty}{\usepackage{xurl}}{} % add URL line breaks if available
\urlstyle{same}
% fallback for those not using the hyperref driver hyperxmp:
\makeatletter
\@ifundefined{xmpquote}{\newcommand{\xmpquote}[1]{#1}}{}
\makeatother
\hypersetup{
  pdftitle={A classification of biologically unavoidable sequences},
  hidelinks,
  pdfcreator={LaTeX via pandoc}}

\title{A classification of biologically unavoidable sequences}
\author{}
\date{10 September 2026}

\begin{document}
\maketitle

\section{Abstract}\label{abstract}

Alexander proved that every eventually periodic sequence occurs in every
infinite gendered population satisfying finite roots, finite outdegree
and finite birthdate sublevels, and asked whether any other sequence has
this property. We construct an avoiding population for each
non-eventually-periodic sequence over every finite alphabet. For a
binary target, the population has two roots, exactly two parents per
non-root and at most two children per vertex. A nonnegative decreasing
path offset proves avoidance. The construction gives a classification by
eventual periodicity. We also characterize universality within the
constructed binary family and deduce finite-observation and
computability obstructions. A Lean development checks the negative
construction and its population axioms; the verification scope is stated
explicitly. No literature priority is claimed.

\section{1. Population axioms and the classification
statement}\label{population-axioms-and-the-classification-statement}

Fix a nonempty finite alphabet \(A\). Following Alexander {[}1,
Definition 1{]}, a population consists of a directed graph \((V,E)\),
edge labels in \(A\), and birthdates \(t:V\to\mathbb R\) such that:

\begin{enumerate}
\def\labelenumi{\arabic{enumi}.}
\tightlist
\item
  There are finitely many vertices with no incoming edge, called roots.
\item
  Every vertex has finitely many children.
\item
  For each real \(R\), the set \(\{v:t(v)\le R\}\) is finite; every edge
  \(u\to v\) satisfies \(t(u)<t(v)\).
\item
  The vertex set is infinite.
\item
  Every non-root has an incoming edge of each label in \(A\).
\end{enumerate}

Each ordered pair of vertices carries at most one edge and one label. A
sequence \(s:\mathbb N\to A\) is realized if there exists a path
\(v_0,v_1,\ldots\) with \(k\)th edge labelled \(s(k)\). There is no
requirement that \(v_0\) be a root. A sequence is biologically
unavoidable when every population over its alphabet realizes it. It is
eventually periodic when there exist integers \(K\ge0\), \(p>0\) such
that \(s(k+p)=s(k)\) for all \(k\ge K\). Throughout this paper,
aperiodic means not eventually periodic.

\textbf{Theorem 1 (classification).} A finite-alphabet sequence is
biologically unavoidable if and only if it is eventually periodic.

The positive direction is Alexander's Theorem 10 {[}1{]}, proved again
in Section 4 with an explicit boundary repair. The negative direction is
the construction in Sections 2--3. It answers the final question of
{[}1, Section 6{]}. The other questions there about universal
populations and ordinal characterizations are not all claimed to be
resolved by this theorem.

\section{2. An explicit binary avoiding
population}\label{an-explicit-binary-avoiding-population}

Fix \(s:\mathbb N\to\{0,1\}\) and put

\[r(2j)=s(j),\qquad r(2j+1)=1-s(j).\]

Let \(P_s\) have vertices \(\mathbb N\) and birthdate \(t(v)=v\). For
every \(v\ge2\) put the following two incoming edges, and no others:

\[v-1\longrightarrow v\quad\text{with label }r(v),\]
\[v-2\longrightarrow v\quad\text{with label }1-r(v).\]

There is no edge \(0\to1\). The roots are exactly 0 and 1. Vertex 0 has
one child, and each vertex \(u\ge1\) has two children. Every non-root
has two distinct parents, with complementary labels. Birthdates increase
strictly, their sublevel sets are finite, and the graph is infinite.
Thus these are populations in the precise edge-labelled sense of
{[}1{]}.

\textbf{Claim.} If \(P_s\) realizes \(s\), then \(s\) is eventually
periodic.

\emph{Proof.} Suppose \(v_0,v_1,\ldots\) is an infinite matching path.
Each increment is 1 or 2. Inductively,

\[v_k\ge 2k.\]

Indeed, if \(v_k\ge2k+1\), the next vertex is at least \(2k+2\). If
\(v_k=2k\), an increment of 1 would have label \(r(2k+1)=1-s(k)\), so is
impossible. At \(k=0\) this edge is absent; the same conclusion holds.
The base case is \(v_0\ge0\).

Consequently \(d_k=v_k-2k\) is a nonnegative integer sequence with
increments in \(\{-1,0\}\). It eventually stabilizes at some \(d\ge0\).
From some \(K\) onward every path increment is therefore 2, so

\[s(k)=1-r(2k+d+2),\qquad k\ge K.\]

If \(d=2e+1\), this says \(s(k)=s(k+e+1)\), giving positive eventual
period \(e+1\). If \(d=2e\), it says \(s(k)=1-s(k+e+1)\); applying the
relation twice gives period \(2e+2\). Either case forces eventual
periodicity.
\leavevmode\unskip\penalty9999\hbox{}\nobreak\hfill\quad\hbox{$\square$}

Thus \textbf{every non-eventually-periodic binary target is avoided by
its own \(P_s\), from all starting vertices}. The graph depends on the
target; this does not produce one graph avoiding every aperiodic
sequence.

\section{3. Finite alphabets and vertex
genders}\label{finite-alphabets-and-vertex-genders}

Let \(A\) be finite and \(s:\mathbb N\to A\) not eventually periodic. At
least one indicator projection \(\pi:A\to\{0,1\}\) has a
non-eventually-periodic image \(\pi\circ s\). Otherwise all finitely
many symbol indicators would be eventually periodic. Beyond the maximum
of their starting thresholds, a common multiple of their positive
periods would be a period of \(s\).

The choice of an aperiodic indicator here is existential. No algorithm
for selecting it from an arbitrary sequence program is asserted. The
binary construction itself is uniform relative to its target; the
computability reductions in Section 5 use only that binary construction.

Construct the binary population \(P_{\pi\circ s}\). Replace each vertex
\(v\) by the copies \((v,a)\), \(a\in A\). For every binary edge
\(u\to v\) with label \(c\), add

\[ (u,a)\longrightarrow(v,b)\quad\text{with label }a
   \quad\text{for all }a,b\in A\text{ satisfying }\pi(a)=c.\]

Give each copy birthdate \(v\). There are \(2|A|\) roots, at most
\(2|A|\) children per vertex, and exactly one incoming parent of each
label at every non-root. The other axioms are inherited. An \(s\)-path
would project to a \(\pi\circ s\)-path in the binary population, which
is impossible.

Moreover, label \(a\) is the source copy's second coordinate: assigning
gender \(a\) to \((v,a)\) makes this a \textbf{vertex-gendered}
population. For a binary alphabet this yields four roots and at most
four children per vertex. The smaller two-root/two-child construction
uses the more general edge-labelled definition.

\textbf{Thue--Morse example.} Put \(t(k)\) equal to the parity of the
number of ones in the binary expansion of \(k\). Then \(t(2k)=t(k)\) and
\(t(2k+1)=1-t(k)\), so the auxiliary sequence \(r\) equals \(t\) itself.
The incoming labels at \(v\) are simply \(t(v)\) and \(1-t(v)\).

For completeness, \(t\) is not eventually periodic. If \(p>0\) were an
eventual period, then for every sufficiently large \(m\),
\(t(2^m-p)=t(2^m)=1\). But \(2^m-p\) is the \(m\)-bit complement of
\(p-1\), so

\[t(2^m-p)=(m\bmod2)\mathbin{\mathrm{xor}}t(p-1),\]

which alternates with \(m\). This is a contradiction.

\section{4. The positive direction and its boundary
repair}\label{the-positive-direction-and-its-boundary-repair}

This section recalls the mathematical input from Alexander {[}1,
Proposition 9 and Theorem 10{]}. It is included to make the
classification argument self-contained and to clarify a deletion step.

Let \(V_j\) consist of vertices reachable in at most \(j\) edges from a
root. Each \(V_j\) is finite by finite roots and finite outdegree. Every
vertex has finitely many ancestors, by the birthdate axiom, and every
vertex descends from a root: a backward path must terminate inside its
finite ancestor set.

\textbf{Proposition 2 (Alexander).} Every periodic sequence is
unavoidable.

\emph{Proof.} Let \(s\) have period \(p>0\) and let
\(w=s(0)\cdots s(p-1)\). Starting at any vertex \(v\), repeatedly try to
select a length-\(p\) path backwards whose forward labels are \(w\).
Whenever a full block can be selected, retain it and repeat from its
first vertex. This process must terminate because all selected vertices
are distinct ancestors of \(v\). If selection stops at \(u\), some
partial backward choice of fewer than \(p\) parents reaches a root, so
\(u\in V_{p-1}\). Otherwise all \(p\) choices would be available by the
incoming-label condition. Thus there is a path from \(V_{p-1}\) to \(v\)
labelled \(w^m\) for some \(m\ge0\).

The lengths of these paths are unbounded as \(v\) varies: a finite set
with finite branching reaches only finitely many vertices within any
bounded number of steps, whereas the population is infinite. The forest
of all \(s\)-matching finite paths starting in the finite set
\(V_{p-1}\) is finitely branching and has unbounded height. König's
lemma gives an infinite \(s\)-matching path. This proves the
proposition.
\leavevmode\unskip\penalty9999\hbox{}\nobreak\hfill\quad\hbox{$\square$}

\textbf{Proposition 3 (finite prefixes).} If \(s\) is unavoidable and
\(w\) is finite, then \(ws\) is unavoidable. Conversely, unavoidability
of \(ws\) implies that of \(s\).

\emph{Proof.} The converse follows by deleting the first \(|w|\) edges
of a realizing path. For the forward implication, take any population
\(P\) and write \(k=|w|\). There is nothing to prove when \(k=0\).
Delete \(V_{k-1}\), then delete every remaining edge entering the finite
boundary \(V_k\backslash V_{k-1}\). Those boundary vertices become
roots. Every vertex outside \(V_k\) retains all its original parents,
since a parent in \(V_{k-1}\) would put its child in \(V_k\). All
population axioms therefore hold in the modified graph \(P'\), which is
infinite because only finitely many vertices were removed.

Realize \(s\) in \(P'\). Its starting vertex lies outside \(V_{k-1}\) in
\(P\). Choose \(k\) parents backwards in \(P\) with labels spelling
\(w\) forwards. No root can be encountered before completing those \(k\)
choices, since that would put the starting vertex in \(V_{k-1}\).
Prepending this path realizes \(ws\) in \(P\).
\leavevmode\unskip\penalty9999\hbox{}\nobreak\hfill\quad\hbox{$\square$}

Applying Proposition 3 to a periodic tail proves the positive direction
of Theorem 1. Merely taking the induced graph after deleting
\(V_{k-1}\), as in the published proof, need not preserve incoming-label
coverage. For example, put two roots \(a_0,b_0\) at time zero and pairs
\(a_i,b_i\) at each later integer time, each normally receiving a
label-0 edge from \(a_{i-1}\) and a label-1 edge from \(b_{i-1}\).
Replace the label-0 parent of \(a_2\) by \(a_0\). After the roots are
deleted, \(a_2\) remains a non-root with no label-0 parent. The
additional boundary-edge deletion above repairs that step without
changing the theorem.

\section{5. Universality and computability in the witness
family}\label{universality-and-computability-in-the-witness-family}

The construction has a second exact property:

\[P_s\text{ realizes every binary sequence}
\quad\Longleftrightarrow\quad s\text{ is eventually periodic}.\]

The negative direction is the binary avoidance theorem in Section 2. If
\(s\) has eventual period \(p\), then \(r\) has eventual period
\(q=2p\). Form a finite residue multigraph modulo \(q\), with the two
incoming transitions into residue \(a\), labelled by the eventual
\(r(a)\) and its complement, and marked with delays 1 and 2. Every state
has an incoming transition of each label. Every finite word can
therefore be read backwards; finite branching and finitely many possible
starts give an infinite path for every infinite word. Starting
sufficiently late in \(P_s\), lift each transition by its marked delay.
This realizes the word in the actual population. The marked delays
retain the information needed to lift the quotient path.

The graph through vertex \(N\) depends only on
\(s(0),\ldots,s(\lfloor N/2\rfloor)\). Any such prefix can be continued
with zeros or with Thue--Morse. The resulting populations agree on that
entire finite region, but one is universal and the other is not. This
also obstructs any finite adaptive sequence of local edge queries.

Even a program describing the graph does not decide universality. Given
a machine \(M\), let \(s_M(k)\) be the \(k\)th Thue--Morse bit if \(M\)
has not halted within \(k\) steps, and zero otherwise. This sequence and
its graph are total computable objects. The graph is universal exactly
when \(M\) halts. Conversely, using zeros until a halt is detected and
then starting a Thue--Morse tail gives a universal graph exactly when
\(M\) never halts. Thus neither universality nor non-universality is
semidecidable on these presentations, even with the promise that all
population axioms hold.

For a fixed aperiodic \(s\), breadth-first matching from any vertex
terminates with an empty set by Section 2 and König's lemma. There is no
bound uniform in the target even for starting vertex 1: a target with
\(N+1\) initial zeros and an aperiodic tail admits \(N\) matching edges
along \(v_k=1+2k\). Obtaining useful explicit bounds for the fixed
Thue--Morse target remains a quantitative question.

\section{6. Verification and
reproducibility}\label{verification-and-reproducibility}

The Lean project defines the actual edge relation and eventual
periodicity, proves the binary path implication, proves the
finite-alphabet indicator reduction, and constructs the lifted
population. The final endpoint is
\texttt{BiologicalAvoidance.finite\_alphabet\_real\_avoidance}, using
explicitly stated real-birthdate population axioms. Consult
\texttt{verification/lean.json} and \texttt{verification/lean.log} for
the check result for this release and exact source hashes.

The formal statement quantifies over every finite alphabet and every
target that is not eventually periodic. It constructs a graph on the
natural numbers times that alphabet, with real birthdates, satisfying
all displayed population conditions and containing no infinite matching
path from any starting vertex. The label predicate is functional on
ordered vertex pairs. The formal proof does not assume an aperiodic
binary projection or a population-existence lemma; those are proved
internally.

Alexander's positive direction, the proof that Thue--Morse is aperiodic,
and the universality/computability corollaries in Section 5 remain prose
proofs. The classification as a whole is not presented as an end-to-end
Lean theorem. \texttt{STATEMENT-AUDIT.md} explains the correspondence
and records that this audit was performed by the same AI system, not an
independent reviewer.

The standard-library Python programs in \texttt{checks/} independently
transcribe edges and matching paths and check finite periodic-template
lifts. These are diagnostic controls, not proofs of infinite claims. For
example, failure to find a barrier within a fixed finite horizon is
recorded as unresolved by that computation, never as an infinite path.

\section{7. Attribution, provenance and
scope}\label{attribution-provenance-and-scope}

Samuel A. Alexander's positive theorem and the original population model
are prior work {[}1{]}. A human collaborator selected the paper and
directed the research. OpenAI Codex (GPT-6 family; no more specific
deployment identifier is asserted) developed the new construction,
arguments, formalization, exact checks and exposition.

Searches described in \texttt{literature-review.md} did not verify an
earlier classification, but do not establish priority. The short
computational reductions are consequences of the construction and the
standard halting obstruction; their inclusion is not a claim that these
methods are novel. No physical or biological prediction is inferred
beyond the population model's stated axioms.

\section{References}\label{references}

{[}1{]} Samuel A. Alexander, \emph{Biologically unavoidable sequences},
Electronic Journal of Combinatorics \textbf{20}(1) (2013), P31.
\href{https://doi.org/10.37236/3035}{DOI: 10.37236/3035}.
\href{https://www.combinatorics.org/ojs/index.php/eljc/article/download/v20i1p31/pdf/}{Published
PDF}. \href{https://arxiv.org/abs/1212.0186}{arXiv:1212.0186}. Theorem
numbers above refer to the published version.

\end{document}
